\documentclass{article}
\usepackage{booktabs}
\usepackage{siunitx}
\usepackage{amsmath}
\usepackage{amssymb}
\usepackage{url}
\usepackage{graphicx}

\begin{document}

\section{COPA Reviewer Response \& Follow-up Experiments}

Below we document additional experiments and comments addressing the reviewer concerns regarding:
1. Comparison with a loose empirical bounds baseline.
2. Comparison with the closest related conformal baseline (\texttt{selfcalibratingconformal}).

\subsection{Clarification to Add to Section 8 of the Paper}

The theoretical contribution of this paper is an interval-valued regression estimator with an auto-calibration guarantee for a Winsorized test label. The experiments have two more limited goals. First, they examine the empirical behaviour of the IVAR intervals, including their width and their containment of the true conditional mean in synthetic settings where this quantity is known. Second, they investigate whether the interval-valued estimates, after minimax merging and cross-fitting, can serve as a useful post-processing method for point prediction.

For each dataset and random seed, we first split the data into an 80\% training portion and a 20\% held-out test portion. All model fitting, calibration, and cross-fitting are performed only on the training portion. Final metrics are evaluated only on the held-out test portion. The uncalibrated baseline labelled ``none'' is trained once on the full training portion. The CVAR methods use the same training portion but internally perform $K = 10$ fold cross-fitting: each fold is used as a calibration set while the remaining folds are used as proper training data. The fold-specific IVAR intervals are merged into point predictions as described in Section 6 and averaged across folds.

Additionally, the reviewer's suggestion to use Conditional Calibration Error (CCE) from van der Laan \& Alaa is not directly applicable to our setting:
\begin{itemize}
    \item Their CCE evaluates coverage calibration of prediction intervals for the future label $Y$.
    \item Our method produces intervals for the conditional mean $E[Y|X]$ (or calibrated selector), not label coverage.
\end{itemize}
Therefore, we do not report CCE as-is because it measures a different notion of validity. The CCE metric evaluates calibration of prediction intervals targeting label coverage, whereas our intervals target the conditional mean. Since the target is not directly observable per test point, coverage-based calibration metrics are not applicable. We instead introduce metrics aligned with our theoretical guarantees.


\subsection{Response Letter Text Additions}

We could also add the following text to our response letter/appendix:
\begin{quote}
We added two compact experiments in response to reviewer comments. First, we compare the proposed Winsorized unbounded construction with a loose-bounds bounded-IVAR baseline, where the empirical calibration-label range is treated as a bound. This tests whether Winsorization provides benefit beyond simply pretending the problem is bounded.

Second, we add a minimal comparison with the public \texttt{selfcalibratingconformal} implementation associated with self-calibrating conformal regression. Because the targets differ, we do not present this as a direct head-to-head replacement comparison. Instead, we report point-prediction metrics where applicable and interval coverage/width for conformal interval outputs. This clarifies the relationship between IVAR/CVAR and the closest available related implementation.
\end{quote}

\subsection{Loose-Bounds Bounded-IVAR Baseline (Part A)}

We compare against a simple bounded-IVAR baseline (\texttt{CVAR-loose-bounds}) that treats the calibration-label range as an empirical bound. In each calibration fold, the lower and upper hypothetical labels are set to the minimum and maximum calibration labels, respectively. This baseline tests whether the unbounded Winsorized construction provides benefit beyond simply pretending that the regression problem is bounded by loose empirical label bounds. The comparison is intentionally limited and is included to clarify the role of Winsorization rather than to provide a complete benchmark of bounded Venn--Abers regression.

The results in Table~\ref{tab:loose_bounds} demonstrate that while the point-prediction RMSE is similar between the methods, the Winsorized construction (\texttt{CVAR1} and \texttt{CVAR10}) yields significantly tighter interval widths than the loose bounds baseline. This suggests that the Winsorized construction improves calibration stability and tightness relative to loose empirical bounds, while retaining the theoretical motivation for unbounded labels.

\begin{table}[htbp]
\centering
\caption{Loose-bounds bounded-IVAR baseline comparison (RMSE and mean interval width, Gradient Boosting base model, $N=10,000$, $\sigma=3$).}
\label{tab:loose_bounds}
\begin{tabular}{lcccc}
\toprule
Dataset & none & CVAR1 & CVAR10 & CVAR-loose-bounds \\
\midrule
Bounded Logistic (RMSE) & 3.168 & 3.125 & 3.129 & 3.125 \\
Bounded Logistic (Width) & — & \textbf{0.146} & \textbf{0.124} & 0.158 \\
\midrule
Linear Gaussian (RMSE) & 3.070 & 3.106 & 3.121 & 3.103 \\
Linear Gaussian (Width) & — & \textbf{0.195} & \textbf{0.157} & 0.215 \\
\midrule
Heavy-tailed (RMSE) & 5.170 & 5.138 & 5.149 & 5.143 \\
Heavy-tailed (Width) & — & \textbf{0.208} & \textbf{0.130} & 0.309 \\
\midrule
Outliers (RMSE) & 5.074 & 5.113 & 5.115 & 5.103 \\
Outliers (Width) & — & \textbf{0.263} & \textbf{0.120} & 0.423 \\
\midrule
Friedman 1 (RMSE) & 3.178 & 3.213 & 3.229 & 3.207 \\
Friedman 1 (Width) & — & \textbf{0.206} & \textbf{0.173} & 0.227 \\
\bottomrule
\end{tabular}
\end{table}

\subsection{Minimal Comparison with Self-Calibrating Conformal Regression (Part B)}

As a minimal comparison with the closest available implementation of self-calibrating conformal regression, we additionally evaluate the public \texttt{selfcalibratingconformal} package (available at \url{https://github.com/Larsvanderlaan/SelfCalibratingConformal}) on three representative synthetic datasets using Gradient Boosting as the shared base model. 

To ensure a fair evaluation, both methods use the same initial 80\% training and 20\% held-out test split, but their internal calibration split protocols differ:
\begin{itemize}
    \item \textbf{CVAR}: Performs $K=10$ fold cross-fitting on the full 80\% training portion. In each fold, 10\% of the data is used as the calibration set while the remaining 90\% serves as proper training data. The predictions are averaged across folds.
    \item \textbf{Self-Calibrating Conformal (SCC)}: Employs a single train/calibration split where the 80\% training portion is split into 75\% proper training data (to train the base Gradient Boosting regressor) and 25\% calibration data (to calibrate the conformal predictor). The target coverage significance level is set to $\alpha = 0.05$ (targeting 95\% coverage).
\end{itemize}

This comparison is intentionally limited because the inferential targets differ. CVAR is evaluated primarily as a merged point regressor derived from IVAR intervals for the Winsorized regression estimate, whereas self-calibrating conformal prediction primarily targets prediction intervals for the future label. We therefore report point-prediction RMSE where applicable and interval coverage/width for conformal interval outputs.

\begin{table}[htbp]
\centering
\caption{Minimal comparison with self-calibrating conformal baseline (means and standard deviations over 100 seeds, $N=10,000$, $\sigma=3$, Gradient Boosting base model).}
\label{tab:scc_comparison}
\resizebox{\textwidth}{!}{
\begin{tabular}{llccccc}
\toprule
Dataset & Method & RMSE($Y$) $\downarrow$ & RMSE($E[Y|X]$) $\downarrow$ & CalibErr $\downarrow$ & Coverage($Y$) & Width \\
\midrule
\textbf{Bounded Logistic} & Base & 3.239 $\pm$ 0.099 & 1.202 $\pm$ 0.235 & 0.675 $\pm$ 0.214 & -- & -- \\
                 & CVAR1 & \textbf{3.144 $\pm$ 0.062} & \textbf{0.933 $\pm$ 0.134} & \textbf{0.229 $\pm$ 0.055} & -- & 0.148 $\pm$ 0.007 \\
                 & SCC-point & 3.158 $\pm$ 0.062 & 0.976 $\pm$ 0.130 & 0.238 $\pm$ 0.057 & -- & -- \\
                 & SCC-interval & 3.179 $\pm$ 0.062 & 1.047 $\pm$ 0.126 & -- & 0.951 $\pm$ 0.007 & 2.726 $\pm$ 0.115 \\
\midrule
\textbf{Heavy-tailed}     & Base & 5.207 $\pm$ 0.772 & 0.954 $\pm$ 0.193 & 0.312 $\pm$ 0.088 & -- & -- \\
                 & CVAR1 & \textbf{5.201 $\pm$ 0.773} & \textbf{0.934 $\pm$ 0.130} & 0.340 $\pm$ 0.087 & -- & 0.201 $\pm$ 0.019 \\
                 & SCC-point & 5.219 $\pm$ 0.772 & 1.029 $\pm$ 0.139 & 0.400 $\pm$ 0.122 & -- & -- \\
                 & SCC-interval & 5.319 $\pm$ 0.766 & 1.422 $\pm$ 0.323 & -- & 0.952 $\pm$ 0.007 & 3.463 $\pm$ 0.249 \\
\midrule
\textbf{Linear Gaussian}  & Base & \textbf{3.073 $\pm$ 0.055} & \textbf{0.664 $\pm$ 0.124} & \textbf{0.229 $\pm$ 0.074} & -- & -- \\
                 & CVAR1 & 3.085 $\pm$ 0.058 & 0.724 $\pm$ 0.136 & 0.249 $\pm$ 0.066 & -- & 0.183 $\pm$ 0.020 \\
                 & SCC-point & 3.091 $\pm$ 0.055 & 0.747 $\pm$ 0.112 & \textbf{0.229 $\pm$ 0.057} & -- & -- \\
                 & SCC-interval & 3.114 $\pm$ 0.058 & 0.838 $\pm$ 0.120 & -- & 0.951 $\pm$ 0.007 & 2.994 $\pm$ 0.295 \\
\bottomrule
\end{tabular}
}
\end{table}

\subsection{Interpretation of Results}

The results in Table~\ref{tab:scc_comparison} should be interpreted as a sanity check rather than a comprehensive benchmark. The self-calibrating conformal method targets prediction intervals for the future label, while IVAR/CVAR targets regression-estimate intervals and their merged point predictions. The comparison confirms that the methods address related but distinct inferential problems.

We observe:
\begin{itemize}
    \item \textbf{Tighter Intervals for Mean Targets}: The interval width of CVAR1 is substantially smaller than that of SCC-interval (e.g. 0.148 vs 2.726). This is expected because CVAR targets the conditional mean $E[Y|X]$ rather than the label $Y$ itself.
    \item \textbf{Point Prediction Performance}: CVAR1 remains highly competitive (and often outperforms both Base and SCC-point) as a point-prediction post-processor despite being derived from a different validity target.
    \item \textbf{Conformal Coverage}: SCC-interval achieves the expected $95\%$ coverage of $Y$, which is consistent with its design target for conformal prediction.
\end{itemize}

\end{document}
